% ============================================================
% intro.tex
% ============================================================

Federal securities class actions represent one of the most
economically significant categories of civil litigation in the
United States, with annual filings fluctuating between 200 and
400 cases and aggregate settlements frequently exceeding \$2
billion \citep{cornerstone2024, nera2024}. 

These cases arise when investors allege violations of federal
securities laws---primarily Section~10(b) of the Securities
Exchange Act of 1934 and SEC Rule 10b-5, which prohibit
fraudulent misrepresentations, or Section~11 of the Securities
Act of 1933, which imposes strict liability for material
misstatements in registration statements. Defendants include
publicly traded corporations and their officers, with plaintiffs
ranging from individual retail investors to large institutional
investors who increasingly serve as lead plaintiffs following the
Private Securities Litigation Reform Act of 1995 (PSLRA)
\citep{cox2008, cheng2010}.

The stakes are substantial. For defendants, direct legal
expenses and settlement payouts range from millions to over a
billion dollars in so-called ``mega-settlements''
\citep{cornerstone2024}. Beyond monetary costs, securities
litigation generates elevated cost of capital, shifts in
investment policies, managerial distraction, and reputational
damage \citep{karpoff2008, graham2008}. For plaintiffs and their
attorneys, duration and outcome directly affect net recovery and
the feasibility of pursuing future cases. For courts and
policymakers, resolution patterns inform debates over the
efficacy of procedural reforms and the balance between deterring
fraud and curbing frivolous litigation.

Despite the prominence of securities class actions, empirical
research on litigation duration and outcomes has been limited in
both methodology and data quality. Practitioner reports provide
valuable descriptive statistics, but do not employ formal survival
analysis or competing-risks frameworks \citep{cornerstone2024,
nera2024}. Academic studies have applied Cox proportional hazards 
models to settlement speed in selected samples \citep{brochet2014}, 
but ignore the competing nature of case dismissal. In practice, 
nearly all core federal securities class actions ultimately either 
settle or are dismissed. Treating these mutually exclusive outcomes
independently introduces selection bias and provides an incomplete 
picture of litigation dynamics. Furthermore, this methodological gap 
is compounded by a profound data challenge: administrative court 
databases frequently misclassify voluntary dismissals, masking 
``hidden settlements'' and distorting true resolution probabilities 
\citep{scales2024}.

This thesis addresses these gaps through four core contributions.
\textbf{First, in methods:} this is, to our knowledge, the first
application of a full competing-risks survival framework to federal
securities class actions using population-level administrative data.
We characterize litigation duration nonparametrically using
Aalen-Johansen cumulative incidence functions (CIFs), estimate
both cause-specific Cox proportional hazards and Fine-Gray
subdistribution models, and move beyond purely associational
estimates through propensity-score composition adjustment and
shared frailty models.

\textbf{Second, in identification:} we apply inverse probability
of treatment weighting (IPTW) to isolate the PSLRA component
from concurrent changes in the observable case mix---the first
such application of propensity-score re-weighting to securities
litigation outcomes. Because the PSLRA is a legislative date
cutoff rather than a treatment influenced by case characteristics,
we interpret the weighted estimates as composition-adjusted
rather than causally identified, and we complement them with
circuit-level shared frailty models that quantify unobserved
jurisdictional heterogeneity.

\textbf{Third, in findings:} we isolate the PSLRA's effect from
over 30~years of non-linear judicial drift, demonstrating
that the Act operates as a persistent accelerant of dismissal
hazard. Across seven robustness specifications---including
alternative disposition coding schemes, temporal restrictions,
circuit sub-samples, and flexible spline time-trend
controls---the PSLRA dismissal hazard ratio ranges from 1.28
to 1.94, with the composition-adjusted IPTW estimate at
approximately 1.53 (a 53\% increase in the instantaneous
rate of dismissal). The PSLRA is also associated with
settlement suppression in the baseline specification (HR $=
0.378$), though this effect attenuates and loses statistical
significance under flexible time-trend controls, suggesting
that part of the raw settlement association reflects secular
changes in litigation patterns. We further quantify how the
PSLRA effect varies over litigation duration and across
circuits, how MDL consolidation bilaterally suppresses both
competing hazards, and how a 1992 placebo cutoff test reveals
pre-existing trends that reinforce our decision to frame all
estimates as associational or composition-adjusted rather than
causally identified.

\textbf{Fourth, in data:} we confront and mitigate the
well-documented challenges inherent to the Federal Judicial
Center's Integrated Database (IDB) \citep{clermont2002,
alexander1991}. We address the ambiguity of disposition codes
through rigorous sensitivity analyses under three alternative
outcome coding schemes and demonstrate that the structural
patterns identified in our survival models are robust to
these underlying administrative measurement uncertainties.